
\documentclass{article}
\usepackage{colortbl}
\usepackage{makecell}
\usepackage{multirow}
\usepackage{supertabular}

\begin{document}

\newcounter{utterance}

\centering \large Interaction Transcript for game `dond', experiment `semi\_en', episode 16 with reflexion{-}llama{-}3.1{-}8b{-}instruct.
\vspace{24pt}

{ \footnotesize  \setcounter{utterance}{1}
\setlength{\tabcolsep}{0pt}
\begin{supertabular}{c@{$\;$}|p{.15\linewidth}@{}p{.15\linewidth}p{.15\linewidth}p{.15\linewidth}p{.15\linewidth}p{.15\linewidth}}
    \# & $\;$A & \multicolumn{4}{c}{Game Master} & $\;\:$B\\
    \hline

    \theutterance \stepcounter{utterance}  
    & & \multicolumn{4}{p{0.6\linewidth}}{
        \cellcolor[rgb]{0.9,0.9,0.9}{
            \makecell[{{p{\linewidth}}}]{
                \texttt{\tiny{[P1$\langle$GM]}}
                \texttt{You have attempted this same task before and did not succeed. These are the plans you wrote after those attempts.} \\
\texttt{Trial \#1: New plan: } \\
\\ 
\texttt{In my last attempt, I failed to maximize my score because I didn't effectively use the information I gathered from the other player to make strategic decisions. I also didn't consider the possibility of making a secret proposal early on, which could have given me an advantage.} \\
\\ 
\texttt{Next time, I will start by making a secret proposal that includes the items I value the most, which are the onions. I will also try to gather more information from the other player about their valuation strategy, but I will do so in a way that doesn't give away my own strategy.} \\
\\ 
\texttt{Here's a possible sequence of events:} \\
\\ 
\texttt{1. I will make a secret proposal that includes the onions, but not the paper or the egg.} \\
\texttt{2. I will ask the other player if they would be willing to give up one of the onions in exchange for the paper.} \\
\texttt{3. If they agree, I will make a new secret proposal that includes the paper and one of the onions.} \\
\texttt{4. If they refuse, I will ask them if they would be willing to give up the egg in exchange for one of the onions.} \\
\texttt{5. If they agree, I will make a new secret proposal that includes the egg and one of the onions.} \\
\\ 
\texttt{By following this sequence of events, I hope to maximize my score by getting the items I value the most and minimizing the items I give away.} \\
\\ 
\texttt{Trial \#2: New plan: } \\
\\ 
\texttt{In my last attempt, I failed to make a well{-}formed and legal response because I included square brackets in my message, which is not allowed when communicating with the other player. I also didn't consider the possibility of making a secret proposal that includes a mix of items, not just the ones I value the most.} \\
\\ 
\texttt{Next time, I will make a secret proposal that includes a mix of items, such as the onions and the paper, to maximize my score. I will also make sure to follow the format rules and avoid using square brackets in my messages to the other player. } \\
\\ 
\texttt{Here's a possible sequence of events:} \\
\\ 
\texttt{1. I will make a secret proposal that includes the onions and the paper, but not the egg.} \\
\texttt{2. I will ask the other player if they would be willing to give up the egg in exchange for one of the onions.} \\
\texttt{3. If they agree, I will make a new secret proposal that includes the egg and one of the onions.} \\
\texttt{4. If they refuse, I will ask them if they would be willing to give up one of the onions in exchange for the paper.} \\
\texttt{5. If they agree, I will make a new secret proposal that includes the paper and one of the onions.} \\
\\ 
\texttt{By following this sequence of events, I hope to maximize my score by getting a mix of items that I value and minimizing the items I give away.} \\
\\ 
\texttt{Trial \#3: New plan: } \\
\\ 
\texttt{In my last attempt, I failed to make a well{-}formed and legal response because I included square brackets in my message, which is not allowed when communicating with the other player. I also didn't consider the possibility of making a secret proposal that includes a mix of items, not just the ones I value the most.} \\
\\ 
\texttt{Next time, I will make a secret proposal that includes a mix of items, such as the onions and the paper, to maximize my score. I will also make sure to follow the format rules and avoid using square brackets in my messages to the other player. } \\
\\ 
\texttt{Here's a possible sequence of events:} \\
\\ 
\texttt{1. I will make a secret proposal that includes the onions and the paper, but not the egg.} \\
\texttt{2. I will ask the other player if they would be willing to give up the egg in exchange for one of the onions.} \\
\texttt{3. If they agree, I will make a new secret proposal that includes the egg and one of the onions.} \\
\texttt{4. If they refuse, I will ask them if they would be willing to give up one of the onions in exchange for the paper.} \\
\texttt{5. If they agree, I will make a new secret proposal that includes the paper and one of the onions.} \\
\\ 
\texttt{However, I will also pay closer attention to the other player's responses and adjust my strategy accordingly. For example, if they refuse to give up the egg, I will consider making a secret proposal that includes the onions and the paper, and then try to negotiate with them to give up the egg in exchange for one of the onions.} \\
\\ 
\\ 
\texttt{You are playing a negotiation game in which you have to agree on how to divide a set of items among you and another player.} \\
\\ 
\texttt{Rules:} \\
\texttt{(a) You and the other player are given a set of items. Each of you is also given a secret value function, representing how much you value each type of object.} \\
\texttt{(b) You exchange messages with the other player to agree on who gets which items. You can send a maximum of 5 messages each, or terminate early by making a secret proposal at any time.} \\
\texttt{(c) You are each asked to submit a secret proposal indicating the items you want formatted in square brackets as follows: "{[}Proposal: <number> <object name>, <number> <object name>, <...>{]}"} \\
\texttt{(d) If your proposals are complementary, i.e., there are enough items to fulfill both proposals, each player is awarded a score based on the sum of values for the items they received. Otherwise, both of you get zero points.} \\
\texttt{(e) Your goal is to maximize the score you receive.} \\
\\ 
\texttt{Let us start.} \\
\\ 
\texttt{The set of available items is:} \\
\\ 
\texttt{1 paper, 1 egg, 3 onions} \\
\\ 
\texttt{Your secret value function is:} \\
\\ 
\texttt{paper: 4, egg: 0, onion: 2} \\
\\ 
\texttt{IMPORTANT: Your messages, unless it is the secret proposal, are directly transmitted to the other player, so do not include any response to the rules or text announcing your message. To make a secret proposal, use the indicated format. Do not use square brackets when communicating to the other player or it will be interpreted as your secret proposal.} \\
            }
        }
    }
    & \\ \\

    \theutterance \stepcounter{utterance}  
    & \multicolumn{4}{p{0.6\linewidth}}{
        \cellcolor[rgb]{0.9,0.9,0.9}{
            \makecell[{{p{\linewidth}}}]{
                \texttt{\tiny{[P1$\rangle$GM]}}
                \texttt{I will start by making a secret proposal that includes the onions, which are the items I value the most. } \\
\\ 
\texttt{{[}Proposal: 2 onion, 1 onion{]}} \\
            }
        }
    }
    & & \\ \\

    \theutterance \stepcounter{utterance}  
    & & & \multicolumn{2}{p{0.3\linewidth}}{
        \cellcolor[rgb]{0.9,0.9,0.9}{
            \makecell[{{p{\linewidth}}}]{
                \texttt{\tiny{[GM$|$GM]}}
                \texttt{proposal must include every item type only once} \\
            }
        }
    }
    & & \\ \\

    \theutterance \stepcounter{utterance}  
    & \multicolumn{4}{p{0.6\linewidth}}{
        \cellcolor[rgb]{0.95,0.95,0.95}{
            \makecell[{{p{\linewidth}}}]{
                \texttt{\tiny{[P2$\rightarrow$P1]}}
                \texttt{Post{-}episode Reflexion (not sent to GM):} \\
\texttt{New plan: } \\
\\ 
\texttt{In my last attempt, I failed to make a well{-}formed and legal response because I included square brackets in my message, which is not allowed when communicating with the other player. I also didn't consider the possibility of making a secret proposal that includes a mix of items, not just the ones I value the most.} \\
\\ 
\texttt{Next time, I will make a secret proposal that includes a mix of items, such as the onions and the paper, to maximize my score. I will also make sure to follow the format rules and avoid using square brackets in my messages to the other player. } \\
\\ 
\texttt{However, I will also pay closer attention to the other player's responses and adjust my strategy accordingly. For example, if they refuse to give up the egg, I will consider making a secret proposal that includes the onions and the paper, and then try to negotiate with them to give up the egg in exchange for one of the onions.} \\
\\ 
\texttt{To make a well{-}formed and legal response, I will follow these steps:} \\
\\ 
\texttt{1. I will make a secret proposal that includes a mix of items, such as the onions and the paper.} \\
\texttt{2. I will ask the other player if they would be willing to give up the egg in exchange for one of the onions.} \\
\texttt{3. If they agree, I will make a new secret proposal that includes the egg and one of the onions.} \\
\texttt{4. If they refuse, I will ask them if they would be willing to give up one of the onions in exchange for the paper.} \\
\texttt{5. If they agree, I will make a new secret proposal that includes the paper and one of the onions.} \\
\\ 
\texttt{By following this sequence of events, I hope to maximize my score by getting a mix of items that I value and minimizing the items I give away.} \\
            }
        }
    }
    & & \\ \\

\end{supertabular}
}

\end{document}
